% Working research draft. Venue cycle and submission year are undecided.
% TODO marks unimplemented methods or experiments, never a measured result.
\documentclass[letterpaper, 10 pt, conference]{ieeeconf}
\IEEEoverridecommandlockouts
\overrideIEEEmargins
\usepackage{graphicx}
\usepackage{amsmath}
\usepackage{amssymb}
\usepackage{booktabs}
\usepackage{xcolor}
\newcommand{\todo}[1]{\textcolor{red}{[TODO: #1]}}
\title{\LARGE \bf DexPrompt: Transferring Dexterous Grasp Knowledge\\
through Learned Demonstration Synthesis to Visual Action Policies}
\author{}
\begin{document}
\maketitle
\thispagestyle{empty}
\pagestyle{empty}

\begin{abstract}
We investigate a two-stage route from an existing dexterous grasp expert to
visually guided grasping of a requested object. The first stage uses exact
simulation geometry and Dexonomy's grasp-type, pose, and contact knowledge to
learn candidate selection and demonstration adaptation. Full approach, closure,
lift, and hold execution determines whether a generated demonstration is accepted.
The second stage trains a DexGraspVLA-style visual action model by behavior
cloning, using cameras, a target-object request, and robot proprioception.
Masks provide an initial goal interface; point and language requests require
an explicit grounding component. The intended evaluation concerns unseen object
geometries and placements in scenes with distractors, without privileged object
geometry at student inference. This is a working research proposal: reproducible
teacher executions and small-scale visual training infrastructure exist, but
learned shape-conditioned selection, demonstration-editing reinforcement learning,
and visual generalization have not been established. We define the comparisons
needed to test whether transferring the expert's knowledge improves demonstration
production and downstream behavior.
\end{abstract}

\section{INTRODUCTION}
An existing grasp expert offers several ways to hold an object, but a feasible
hand pose alone does not specify how an arm should approach it or whether the
object remains held during lifting. Furthermore, a policy acting on camera
observations must identify the requested object and execute the motion despite
occlusion and tracking error. These are two related learning problems rather
than a single pose-prediction task.

DexPrompt proposes to use privileged simulation during data generation and
visual observations during deployment. Stage 1 learns which Dexonomy grasp
candidate to use and how to adapt its demonstration to the shape and scene.
Stage 2 learns continuous arm and hand actions from the resulting demonstrations.
The target embodiment is the FFW-SH5 / HX5-D20. The initial instruction selects
an object; part masks, specified contact regions, and downstream tool use are
optional extensions, not prerequisites for this study.

Our hypothesis is that explicit type/contact knowledge improves the yield,
diversity, or cost of executable demonstrations, and that these improvements
benefit a visual student under controlled training budgets. Neither using
Dexonomy nor adding reinforcement learning (RL) establishes this benefit by
itself. Prior work already studies grasp-conditioned execution, demonstration
editing, and visual distillation~\cite{unidexgrasp,demofungrasp,clutterdexgrasp}.
Dexonomy-based pretraining for manipulation is also a direct
precedent~\cite{fromgraspstodexterity}.

The planned contributions are a shape-conditioned selection/adaptation mechanism,
a traceable execution-to-visual-data pipeline, and controlled evidence of the
knowledge transferred. These remain objectives pending experiments. The phrase
``unseen objects'' will refer to a declared geometry split and a bounded robot
workspace, not universal physical graspability.

\section{RELATED WORK}
\textbf{Taxonomy and candidate generation.}
Dexonomy provides taxonomy-based synthesis and studies type classification and
type-conditioned grasp generation~\cite{dexonomy}. Its learned components are a
reference for selecting grasp types from shape, but a type score or generated
pose is not a calibrated probability of full arm-hand execution success.
DexGraspNet 2.0 provides scene-conditioned generative grasping~\cite{dexgraspnet2}.
Such generators can supply candidates; motion execution and validation remain
separate stages in our proposed comparison.

\textbf{Demonstration adaptation and privileged execution.}
DemoFunGrasp edits demonstrations with one-step PPO~\cite{demofungrasp}. Here
``one-step'' means selecting an edit and evaluating its executed trajectory,
not completing a grasp in one simulator tick. Grasp style is supplied to that
method; automatically choosing a Dexonomy type is a separate problem in ours.
Its editing formulation is a principal reference for Stage 1, not merely
background supporting the use of simulation.

ClutterDexGrasp learns a privileged RL teacher and distills it into a visual
student~\cite{clutterdexgrasp}. Its shape-aware observations and curriculum
motivate possible closed-loop extensions. It does not use Dexonomy as its
teacher. Differences in embodiment, physics, and contact treatment prevent
assuming that its settings transfer unchanged to our executor.

\textbf{Visual action learning.}
DexGraspVLA combines a language/vision planner, object masks, DINOv2 features,
and a diffusion action controller trained from demonstrations~\cite{dexgraspvla}.
Thus mask and DINO features are consistent with a VLA-style architecture;
using them does not create a competing alternative to VLA. GraspVLA provides a
reference for synthetic action pretraining~\cite{graspvla}. DexVLG predicts
language-conditioned grasp poses~\cite{dexvlg}; pose prediction and full
closed-loop execution require different metrics.
From Grasps to Dexterity uses pretrained grasp representations and hierarchical
keypoint-conditioned behavior learning~\cite{fromgraspstodexterity}. That is a
concrete alternative architecture, not evidence that our current student already
implements its method or inherits its generalization.

\section{TASK AND INFORMATION CONTRACT}
Let $\mathcal{M}$ be exact object geometry, $T_o$ its simulated pose, $C$ the
scene, and $q_t\in\mathbb{R}^{27}$ the arm and finger joint configuration.
The synthesis stage may use all these variables, contact measurements, and
Dexonomy templates. It records both successful executions and failures.
The visual student receives camera images $I_t$, proprioception $q_t$, and a
request $u$. A grounding module produces a target-object mask $M_t$ from a
provided mask, an image point, or language. A user-selected image point is not
an object point cloud and does not impose a contact location.

The student emits arm/hand actions rather than just a final grasp pose.
Privileged mesh, object pose, teacher targets, future phase labels, and future
contact states are unavailable at normal student inference. Their use in a
diagnostic must be labeled as oracle input. Ground-truth simulator masks can
provide supervision and an oracle-goal baseline, but predicted grounding must
be evaluated separately before claiming autonomous point/language operation.

The primary event is selecting the correct object and completing the defined
approach--grasp--lift--hold task under the frozen quality contract. Required
teacher contact groups remain part of the collection contract and a separate
transfer metric. They do not certify an exact taxonomy label or mathematical
force closure. Requesting a watering can alone does not require its handle;
that would need a separately specified semantic or part constraint.

\section{STAGE 1: LEARNED DEMONSTRATION SYNTHESIS}
\subsection{Candidate Selection}
A Dexonomy-derived candidate pool $\mathcal{G}$ contains type, wrist pose,
pregrasp/grasp/squeeze configurations, and required contact structure. The
proposed selector is
\begin{equation}
g \sim s_\phi(g\mid\mathcal{M},T_o,q_0,C,\mathcal{G}).
\end{equation}
A nonlearned ranking and a learned ranking must receive identical candidate
pools and dynamic execution budgets. Full execution outcomes provide labels;
passing a static refiner alone does not. Only attempted candidates have
observed dynamic labels, so selection bias and unexecuted candidates must be
reported rather than treating all untried candidates as failures.
\todo{Specify geometry encoding, type/candidate heads, training labels, splits,
and a nonlearned ranking baseline before fitting the selector.}

\subsection{Demonstration Editing}
The proposed adaptation policy selects bounded edits to a candidate's reference
motion, after which a feedback executor evaluates the entire trajectory:
\begin{align}
\delta &\sim \pi_\eta(\delta\mid\mathcal{M},T_o,q_0,C,g,\tau_g),\\
\tau' &= E(\tau_g,\delta;C).
\end{align}
This is inspired by demonstration-editing RL~\cite{demofungrasp}. The edit
parameterization is unresolved; pose, timing, and joint edits must respect
kinematics and preserve an explicit relationship to the selected teacher.
The reward must reflect full execution success, quality, and cost. Pose
closeness cannot replace held-object success.
\todo{Freeze edit variables, bounds, observations, reward, and rollout budget.
Validate the unedited action reproduces the baseline before training.}

Type selection and adaptation have separate ablations. First compare rankings;
then compare unedited and edited executions of the same selected candidates.
A fixed-candidate, seven-joint approach policy is a restricted execution tool,
not an implementation or test of this full shape-conditioned synthesis stage.

\subsection{Execution and Data Acceptance}
The existing executor checks approach feasibility and uses contact feedback
during closure and lift. It is retained as a baseline. Teacher squeeze joint
angles do not replace contact feedback: a stored target need not remain
appropriate after the object moves in the hand. No controller or acceptance
threshold is relaxed to increase the reported yield.

Each attempt retains object/mesh/scale/placement, type/template, parent candidate,
edit, model and initial-state hashes, physics settings, phase outcomes, contact
signals, and quality measurements. Reproduction checks are reported separately
from new attempts. Filtering, initialization padding, refinement, execution,
and acceptance have distinct counts. Different contact-physics configurations
remain separate datasets until an explicit comparison justifies their use.

\section{STAGE 2: VISUAL BEHAVIOR CLONING}
The visual policy predicts action chunks from available observations:
\begin{equation}
\hat a_{t:t+H-1}=f_\theta(I_{\leq t},M_{\leq t},q_{\leq t}).
\end{equation}
Behavior cloning (BC) is the learning objective family; a diffusion action head
is one model implementation. RGB, masks, and pretrained DINO features are
compatible with the DexGraspVLA reference design. Our current implementation
combines a pretrained visual encoder with a newly trained action model; it does
not transfer the original VLA's complete pretrained action checkpoint.

Accepted teacher executions undergo separate student intake, observation
rendering, action/state alignment, and visibility checks. Temporary goal
occlusion must be handled using available observation history, without revealing
future simulator state. Existing optional part-condition fields may remain for
compatibility, but object-only data and evaluation do not require them.

Low training loss or finite full-model outputs are integration checks, not
closed-loop grasp results. Evaluation must execute the learned actions, measure
correct-target lift and hold, and report distribution shift from demonstrations.
A fixed-object diagnostic is useful for debugging this path but cannot stand in
for Stage 1 selection learning or unseen-object evaluation.
\todo{Freeze architecture, loss parameterization, action decoding, grounding,
and dataset split for each experiment; retain sampler and checkpoint metadata.}

\section{CURRENT EVIDENCE AND EXPERIMENT PLAN}
\subsection{Development Evidence, Not Final Results}
As recorded in the development reports at this draft revision, a small HOPE A
student-intake pool contains 12 demonstrations, 12,477 frames, three objects,
and five teacher candidates. This is a diagnostic training pool, not an
independent generalization benchmark. Reproducible full teacher executions have
been obtained for some objects; static acceptance remains distinct from
execution success and student readiness.

Prior fixed-task approach experiments transferred a scripted approach with BC;
PPO gains depended on training seed. They did not train the proposed
shape-conditioned selector or editing policy. Visual-model training and
evaluation infrastructure exists, but unseen-object visual success is not
established. Ongoing-job outcomes must be read from their final archived
reports before adding result tables. A separate Milk experiment uses different
contact physics and fails a sustained initial-stability check; teacher success
there does not authorize merging it into the A student pool.

\subsection{Controlled Comparisons}
\begin{enumerate}
\item Same pool and execution budget: nonlearned versus learned selection.
\item Same candidates and physics: unedited versus learned demonstration edits.
\item Same student and training budget: data produced by the compared synthesis
methods, reporting both equal-data and collection-cost effects where feasible.
\item Type/contact-informed transfer versus a clearly defined reduced-knowledge
baseline, without silently changing acceptance criteria.
\item Oracle target masks versus predicted point/language grounding; change only
the requested object in an otherwise identical scene to test goal dependence.
\end{enumerate}
The main outcome joins correct-target selection, lift, and hold. Report contact
compliance, quality margins, wrong-object interactions, failures by phase,
attempts, simulator steps, and wall-clock/compute costs separately.

Split source geometries before generation; a new scale, texture, or view of the
same mesh is not a new held-out geometry. Separate training, development/model
selection, and final evaluation. Fresh action-noise seeds on a training object
measure stochastic robustness rather than object generalization. Use multiple
training seeds and uncertainty estimates. A checkpoint-selection procedure is
part of the evaluated method, including fallback to its BC baseline.
\todo{Freeze unseen geometries, layouts, physical feasibility scope, budgets,
metrics, and model-selection rules before final evaluation.}

\subsection{Deployment}
Begin with object-selection tasks and modest clutter, then evaluate occlusion,
unseen shapes, and real-robot transfer. Point and language performance includes
the grounding component; supplied oracle masks cannot certify it. Hardware
calibration, sensing, and any real-data adaptation must be disclosed separately.
\todo{Establish simulator closed-loop results before claiming hardware transfer.}

\section{LIMITATIONS AND OPEN DECISIONS}
Learned synthesis and visual generalization are both unfinished. Current data
cover few shapes and selected successful candidates. Grasp templates, contact
models, object scale/mass assumptions, and acceptance gates can restrict yield.
Reproducible simulation does not prove physical realism. A well-tracking hand
may still lose its object, and contact presence alone does not imply robust
load support. External paper results do not guarantee transfer to this robot.

The intended contribution must exceed assembling existing modules: controlled
experiments must identify what Dexonomy knowledge improves and at what cost.
If editing RL adds no benefit, that result should be reported instead of
presenting RL as intrinsically necessary. Part requests and functional tool use
are optional future scope, not missing prerequisites for the core experiment.

\section{CONCLUSION}
DexPrompt proposes learned shape-conditioned selection and adaptation of
Dexonomy grasps for automatic demonstration production, followed by visual BC
for requested-object grasping. The two learning stages and their evidence must
be assessed separately and then connected by controlled downstream evaluation.
\todo{Replace this proposal conclusion with measured findings after evaluation.}

\begin{thebibliography}{99}
\bibitem{dexonomy} J. Chen \emph{et al.},
\emph{Dexonomy: Synthesizing All Dexterous Grasp Types in a Grasp Taxonomy},
RSS, 2025. arXiv:2504.18829.
\bibitem{dexgraspvla} Y. Zhong \emph{et al.},
\emph{DexGraspVLA: A Vision-Language-Action Framework Towards General Dexterous
Grasping}, arXiv:2502.20900, 2025.
\bibitem{clutterdexgrasp} Z. Chen \emph{et al.},
\emph{ClutterDexGrasp: A Sim-to-Real System for General Dexterous Grasping in
Cluttered Scenes}, CoRL, 2025. arXiv:2506.14317.
\bibitem{demofungrasp} C. Mao \emph{et al.},
\emph{Universal Dexterous Functional Grasping via Demonstration-Editing
Reinforcement Learning}, arXiv:2512.13380, 2025.
\bibitem{unidexgrasp} Y. Xu \emph{et al.},
\emph{UniDexGrasp: Universal Robotic Dexterous Grasping via Learning Diverse
Proposal Generation and Goal-Conditioned Policy}, CVPR, 2023. arXiv:2303.00938.
\bibitem{dexvlg} J. He \emph{et al.},
\emph{DexVLG: Dexterous Vision-Language-Grasp Model at Scale},
arXiv:2507.02747, 2025.
\bibitem{graspvla} S. Deng \emph{et al.},
\emph{GraspVLA: a Grasping Foundation Model Pre-trained on Billion-scale
Synthetic Action Data}, arXiv:2505.03233, 2025.

\bibitem{fromgraspstodexterity} Y. Yuan, X. Liu, S. Krishna, and D. Held,
\emph{From Grasps to Dexterity: Large-Scale Grasp
Pretraining for Dexterous Manipulation}, arXiv:2606.30749, 2026.
\bibitem{dexgraspnet2} J. Zhang \emph{et al.},
\emph{DexGraspNet 2.0: Learning Generative Dexterous Grasping in Large-scale
Synthetic Cluttered Scenes}, CoRL, 2024. arXiv:2410.23004.
\end{thebibliography}
\end{document}
